\addtocontents{toc}{\protect}
\section{\MakeUppercase{conclusions}}
\label{sec:chapter-conclusions}

\counterwithin{figure}{section}
\counterwithin{table}{section}


\subsection{Black Hole Mergers in AGN Disks}

We used the new public, open-source, reproducible code \code{McFACTS} to test the $q$--$\Xeff$ distribution for the LVK AGN channel.
Broadly we find islands of structure in $q$--$\Xeff$ parameter space as a function of BH generation, spin, eccentricity distributions, disk models, and AGN lifetime.
The observed $q$--$\Xeff$ anti-correlation can be well reproduced by a dense and relatively short-lived model AGN disk with some dynamical cooling.
Tighter constraints on the $q$--$\Xeff$ distribution from future LVK BBH catalogs (O4 and O5) will allow us to strongly constrain key AGN disk and NSC properties.

\subsection{Supernovae in AGN Disks}

Massive stars in nuclear star clusters on orbits embedded in AGN accretion disks may detonate as SN from 
within the disk.
The shock can break free from the disk if its internal pressure exceeds the ambient disk pressure and is muffled otherwise.
The exact outcome depends on the \textbf{radial and vertical location} of the detonation as described by Eq.\,(\ref{eqn:C-parameterized}).
We performed hydrodynamic three-dimensional shearing box simulations at two locations on either side of the muffling boundary of embedded SNe and found agreement with the analytic prediction for shock breakout from an SG-type AGN disk orbiting a black hole of mass $\Mbh=10^8\,\Msun$.
Specifically, we found:

\begin{itemize}
    \item In general, SNe break free from the inner disk, whereas the outskirts of AGN disks muffle them.
    \item In SG disks, when $\mbh=10^6\,\msun$, the muffling boundary occurs at $R_{\rm m}=5\times10^6\,\rs$ ($\sim10^6$ AU) and moves inward to $R_{\rm m}=100\,\Rs$ ($\sim100$ AU) in disks around $\mbh=10^9\,\msun$.
    \item In TQM disks, the muffling radius lies at $R_{\rm m}=10^6\,\Rs$, regardless of the SMBH's mass.
    \item The muffling radius for SNe that occur at $z>1\,H$ shifts outwards by approximately one order of magnitude in comparison with midplane detonations.
    \item \textbf{When SNe detonate away from the midplane the ejecta traveling to the far side of the disk may remain shocked or emerge as subsonic flows depending on the gas density and scale height, where the optical thickness will determine whether photons do the same.}
    \item When the SN occurs away from the midplane, the shock propagating away from the disk expands rapidly in comparison to the shock propagating towards (and possibly through) the midplane, slowing down as it climbs the density and pressure gradients.
    \item Considering a lifetime of $\tau_{\rm AGN}=1$ Myr and a SN rate in the disk of $\dot{N}_{\rm SN} \sim 10^{-2}$, we estimate a rate of $\mathcal{R}_{\rm b} \sim 10$ SN AGN$^{-1}$ yr$^{-1}$ breaking free from the disk and estimate a volumetric rate density of $\mathcal{R}_{\rm SN} \sim 10$ Gpc$^{-3}$ yr$^{-1}$ for the local Universe.
\end{itemize}

For SNe at the midplane in the inner disk ($r=10^3\,Rs$), the radiation emerges $\sim30$ days later and rises quickly to a peak at $L=10^{46}\ \text{erg s}^{-1}$.
Typical AGN luminosities can be $\mathcal{O}(10^{43}-10^{48}\ \text{erg s}^{-1})$, so these SNe will only be brighter than AGN of $L\ll10^{46}$, which is the range for Seyfert galaxies and other low-luminosity AGN, and will not be visible in quasars.
Radiation from a SN at the same location but offset from the midplane at $z=1\,H$ emerges 15 days earlier, rises much faster, and peaks at $2 \times 10^{46}\ \text{erg s}^{-1}$ when viewed from the same side of the midplane.
The radiation does not reach the other side of the disk in a detectable manner under this scenario.
While bright, these SNe may not be detectable against a bright AGN disk, but would be easily seen if occurring in a starved disk or during an episode of lower mass accretion rates.
For SNe detonating in the outer disk at $r=10^5\,\Rs$, the large scale height of the disk prevents the hot shock from emerging from the photosphere.

Embedded SNe do not exhibit a significant color evolution and cause the observed patch to become bluer rather than redder, as is typical of SNe, upon emerging from the disk.
Without this key indicator that is common to naked SNe, they may be indistinguishable from flares produced by recoiling BBH merger remnants.
As such we will need to consider a different part of the spectrum or aspects of these transients to reliably distinguish between flares emitted by embedded merger remnants versus embedded SNe in transient searches and archives such as ZTF or LSST.
